\GalSourcePageBoundary{063}{L0062}{34}{34}

équations dont le degré est un nombre premier. Voici le théorème
donné par notre analyse:

\begin{quote}\itshape
Pour qu'une équation de degré premier, qui n'a pas de diviseurs
commensurables, soit soluble par radicaux, il faut et
il suffit que toutes les racines soient des fonctions rationnelles
de deux quelconques d'entre elles.
\end{quote}

Les autres applications de la théorie sont elles-mêmes autant
de théories particulières. Elles nécessitent d'ailleurs l'emploi de
la théorie des nombres, et d'un algorithme particulier: nous les
réservons pour une autre occasion. Elles sont en partie relatives
aux équations modulaires de la théorie des fonctions elliptiques,
que nous démontrons ne pouvoir se résoudre par radicaux.

\medskip
\noindent\hspace{3em}Ce 16 janvier 1831.\hfill \textsc{E. Galois.}\hspace{3em}

\begin{center}
\rule{0.15\linewidth}{0.4pt}\par
\vspace{0.7em}
{\bfseries PRINCIPES.}
\end{center}

Je commencerai par établir quelques définitions et une suite de
lemmes qui sont tous connus.

\emph{Définitions.} --- Une équation est dite \emph{réductible} quand elle
admet des diviseurs rationnels; \emph{irréductible} dans le cas contraire.

Il faut ici expliquer ce qu'on doit entendre par le mot \emph{rationnel},
car il se représentera souvent.

Quand l'équation a \emph{tous} ses coefficients numériques et rationnels,
cela veut dire simplement que l'équation peut se décomposer
en facteurs qui aient leurs coefficients numériques et rationnels.

Mais quand les coefficients d'une équation ne seront \emph{pas tous}
numériques et rationnels, alors il faudra entendre par diviseur
rationnel un diviseur dont les coefficients s'exprimeraient en
fonction rationnelle des coefficients de la proposée.\GalPrintedOmission{W09-WV001}

Il y a plus: on pourra convenir de regarder comme rationnelle
toute fonction rationnelle d'un certain nombre de quantités déterminées,
supposées connues \emph{a priori}. Par exemple, on pourra
